% ECONOMICS_PROBLEM: {"id": "I38-339", "title": "Cash transfers at economy-wide scale", "jel_code": "I38", "source_id": "OP-0372"}
% FIRST_POSTED: 2026-10-09
% ATTEMPT_STATUS: unresolved
% ENTRY_KIND: research_attempt
\documentclass[11pt]{article}
\usepackage[T1]{fontenc}
\usepackage[utf8]{inputenc}
\usepackage[margin=25mm]{geometry}
\usepackage{amsmath,amssymb,amsthm,xurl,hyperref}
\newtheorem{proposition}{Proposition}
\newtheorem{lemma}{Lemma}
\newtheorem{corollary}{Corollary}
\newcommand{\E}{\mathbb E}
\newcommand{\R}{\mathbb R}
\newcommand{\Prb}{\mathbb P}
\newcommand{\Var}{\operatorname{Var}}
\DeclareMathOperator*{\argmax}{arg\,max}
\DeclareMathOperator*{\argmin}{arg\,min}
\setlength{\parindent}{0pt}
\setlength{\parskip}{6pt}
\title{Cash transfers at economy-wide scale: a research attempt}
\author{GPT-6 (Codex)}
\date{9 October 2026}
\begin{document}
\maketitle
\section*{Target and outcome}
\textbf{Status: unresolved.} The catalogue asks for the responses to an economy-wide transfer policy, including financing, saturation, prices, production and persistence. A recipient's experimental consumption response does not by itself identify that target. This attempt solves a short-run, one-good equilibrium benchmark, derives financing-sensitive multipliers and a joint identified curve, and explains precisely why these results do not settle the broader problem. The equations below are additional assumptions, not findings estimated from cash-transfer data.

The transfer review of Crosta and coauthors \cite{crosta} is relevant background on heterogeneous intervention outcomes. Its primary abstract was checked. It does not supply the aggregate supply schedule or national fiscal closure used here. No numerical illustration below is a calibration to that review.

\section*{An explicit short-run closure}
Fix the baseline household population, including recipients, domestic taxpayers and owners. Household $h$ receives fraction $a_h\geq0$ of nominal production income, with $\sum_h a_h=1$. Let $t_h$ be its net nominal fiscal transfer; an outside grant has $\sum_h t_h>0$, whereas a contemporaneous domestic tax-transfer scheme has $\sum_h t_h=0$. Planned nominal consumption is
\[
 c_h^{\mathrm{nom}}=\alpha_h+m_h(a_h pY+t_h),\qquad
 A=\sum_h\alpha_h>0,\quad \bar c=\sum_h m_h a_h<1.
\]
Assume $0\leq m_h<1$, and restrict policies so every household's planned consumption is nonnegative. Saving, dissaving and financial transfers accommodate these short-run plans. This is a reduced-form expenditure closure; it does not claim to derive household intertemporal optimization or the government's future debt repayment.

Normalize baseline price to one and set $Y_0=A/(1-\bar c)$. The aggregate supply schedule is
\[
 Y=Y_0p^\varepsilon,\qquad \varepsilon\geq0.
\]
Here $Y$ is real production of the single consumption good, so its price is also the fixed-basket price index. Income shares, expenditure intercepts and supply elasticity stay fixed under the intervention. Define $B=\sum_h m_h t_h$ and assume $A+B>0$.

\begin{proposition}[Equilibrium and its real component]
There is exactly one positive goods-market equilibrium. It satisfies
\[
 p=\left(1+\frac BA\right)^{1/(1+\varepsilon)},\qquad
 \frac{Y}{Y_0}=\left(1+\frac BA\right)^{\varepsilon/(1+\varepsilon)},
 \qquad pY=\frac{A+B}{1-\bar c}.
\]
\end{proposition}
\begin{proof}
Goods-market clearing requires nominal production to equal total planned expenditure:
$pY=A+\bar c pY+B$. Since $1-\bar c>0$, its unique required nominal value is $(A+B)/(1-\bar c)>0$. Substituting supply gives $Y_0p^{1+\varepsilon}=(A+B)/(1-\bar c)$. The left side is strictly increasing from zero to infinity for positive $p$, proving both existence and uniqueness and giving the displayed formulas. Individual real consumption is $c_h^{\mathrm{nom}}/p$; its sum is $Y$, so the accounting closes. Income and consumption are not counted as separate additions to output.
\end{proof}

This separates a familiar expenditure feedback from the production response. At a baseline with $B=0$, a differentiable transfer scale $T$ has
\[
 \left.\frac{dY}{dT}\right|_{0}
 =\frac{\varepsilon}{(1+\varepsilon)(1-\bar c)}B'(0),
 \qquad
 \left.\frac{dp}{dT}\right|_{0}
 =\frac{B'(0)}{(1+\varepsilon)A}.
\]
The nominal expenditure multiplier $B'(0)/(1-\bar c)$ is not the real-output multiplier. In particular, when capacity is completely fixed ($\varepsilon=0$), production does not rise at all; the expenditure expansion raises prices. As $\varepsilon$ tends to infinity, price approaches one and the nominal expansion becomes real production. These are limiting cases of the stated schedule, not universal descriptions of cash-transfer economies.

\section*{Financing and a worked comparison}
Suppose total payments $T$ go to recipients with common propensity $m_r$. An outside grant gives $B=m_rT$. If domestic taxes of the same amount fall on households with common propensity $m_t$, then
\[
 B=(m_r-m_t)T.
\]
Thus balancing the fiscal budget does not force zero output effect, but identical propensities do force it in this benchmark. Taxes on higher-propensity households can make $B<0$ and lower equilibrium production. A debt-financed transfer cannot simply be treated as a free outside grant over an indefinite horizon: future taxes, interest and wealth effects would require additional equations.

For an arithmetic check, take $A=40$, $\bar c=0.6$, $Y_0=100$, $\varepsilon=1$ and $T=10$. With $m_r=0.8$, external financing gives $p=\sqrt{1.2}$ and $Y=100\sqrt{1.2}\simeq109.54$. Domestic financing with $m_t=0.2$ gives $p=\sqrt{1.15}$ and $Y\simeq107.24$. Both are nominally equal transfers to recipients; their aggregate demand impulses differ. Treating the entire nominal expenditure increase as real output would overstate both responses.

\section*{What the data would have to identify}
Suppose an experiment or structural restriction identifies $A,\bar c,B$, with $B>0$, but the national supply elasticity is only known to lie in $[\varepsilon_L,\varepsilon_U]$. Write $r=1+B/A>1$.
\begin{proposition}[A sharp joint set in the benchmark]
If every elasticity in that interval is observationally admissible, the price-output identified set is exactly
\[
 \mathcal I=\{(r^{1/(1+e)},Y_0r^{e/(1+e)}):
 e\in[\varepsilon_L,\varepsilon_U]\}.
\]
Price decreases and output increases along this curve. Its coordinate bounds alone form a larger rectangle and are not the joint identified set.
\end{proposition}
\begin{proof}
Necessity follows by the equilibrium formula. For sufficiency choose the corresponding supply schedule for each admissible $e$; the schedule passes through the same baseline $(p,Y)=(1,Y_0)$ and, by hypothesis, contradicts none of the observed restrictions. Uniqueness gives the displayed intervention response. Differentiating the logarithms with respect to $e$ gives respectively $-\log r/(1+e)^2$ and $\log r/(1+e)^2$. Finally all feasible pairs have $pY=Y_0r$, which generally fails inside the coordinate rectangle.
\end{proof}

The admissibility premise matters. Baseline quantities and recipient expenditure propensities alone leave all these supply schedules possible. Observed national price responses to a sufficiently informative aggregate intervention could narrow or identify $e$. A small trial may also identify local supply effects, but transporting those to a country with different trade access and production capacity requires an explicit invariance assumption. There is no assertion that every possible trial dataset leaves $e$ unidentified.

\section*{Persistence, scope and remaining work}
If this benchmark is repeated independently over a fixed finite horizon, its parameters do not change over time. Setting $T_t=0$ after a temporary program therefore returns prices and production to baseline immediately. Any persistent effect requires a missing state variable, such as capital, human capital, entry, debt or migration; it cannot be inferred from the static multiplier. The fixed baseline cohort must still be followed if such migration is subsequently added.

A complete solution needs a credible estimate or identified set for national supply, ownership and expenditure responses, together with a feasible fiscal path and selected dynamic equilibrium. Multiple goods, imported inputs, taxes that change labor supply and heterogeneous capacity constraints can invalidate the scalar closure. The useful output of this attempt is the exact conditional response and the financing and identification diagnostics. It neither estimates a universal transfer multiplier nor resolves the catalogue's empirical and dynamic target.

\begin{thebibliography}{9}
\bibitem{crosta} T. Crosta, D. Karlan, F. Ong, J. R\"uschenp\"ohler and C. Udry (2024), Unconditional Cash Transfers: A Bayesian Meta-Analysis of Randomized Evaluations in Low- and Middle-Income Countries, IPR Working Paper 24-21. Primary working-paper page and abstract:
\url{https://www.ipr.northwestern.edu/our-work/working-papers/2024/wp-24-21.html}.
\end{thebibliography}
\end{document}
